### Literal statement (as written)

Let (A={a_1<a_2<\cdots}) be an infinite strictly increasing sequence of **integers**.
For each integer (n), let (f(n)) be the number of pairs ((u,v)) of indices with (1\le u\le v) such that
[
n=\sum_{u\le i\le v} a_i.
]
Question: does there exist such an (A) for which (f(n)\to\infty) as (n\to\infty)?
Or even (f(n)\ge 2) for all sufficiently large (n)?

### Ambiguities / stress points

1. **Sign of the (a_i).** The sequence is said to be “integers,” not “positive integers.” If negative terms are allowed, then for large positive (n) only finitely many negative terms can be used (since the sequence is strictly increasing, it contains only finitely many negative integers), so the asymptotic question is essentially about the positive tail anyway. But many standard reformulations (convex sets, convolutions on (\mathbb Z_{\ge0})) implicitly assume positivity.
2. **Domain of (n).** The question “as (n\to\infty)” indicates (n\in\mathbb N) (positive integers). I will take (n\ge 1).
3. **Does (u=v) count?** Yes: the definition allows (u=v). This matches the example (a_n=n), where (f(n)) equals the number of odd divisors of (n), which includes the divisor (1) (corresponding to the length-(1) representation).

### Minimal corrected statement (standard conventions)

I will work with:

> **Corrected problem.**
> Let ((a_i)_{i\ge1}) be a strictly increasing sequence of **positive** integers. For each (n\in\mathbb N), define
> [
> f(n):=\#\{(u,v)\in\mathbb N^2: 1\le u\le v,\ \sum_{i=u}^v a_i = n\}.
> ]
> Does there exist such ((a_i)) with (f(n)\to\infty) as (n\to\infty)?
> Or even with (f(n)\ge 2) for all sufficiently large (n)?

This is the interpretation consistent with the “convex set” reformulation and with all cited examples.

### Equivalent convex/difference formulation (made explicit)

Define prefix sums
[
s_0:=0,\qquad s_m:=\sum_{i=1}^m a_i\quad(m\ge1),
]
and let (S:={s_m:m\ge0}\subset\mathbb Z).

Then for (1\le u\le v),
[
\sum_{i=u}^v a_i = s_v-s_{u-1}.
]
Hence
[
f(n)=\#\{(i,j)\in\mathbb Z_{\ge0}^2: i<j,\ s_j-s_i=n\}.
]
In additive combinatorics notation, if (1_S) is the indicator of (S), then
[
(1_S\circ 1_S)(n):=\sum_x 1_S(x),1_S(x+n) = \#\{(x,y)\in S^2: y-x=n\},
]
so (f(n)=(1_S\circ 1_S)(n)) (up to the harmless convention of ordered vs. (i<j)).

Moreover, because (a_{m+1}>a_m), the consecutive gaps (s_{m+1}-s_m=a_{m+1}) are strictly increasing, so (S) is a **convex** set in the standard discrete sense (“successive differences strictly increasing”).

So the corrected problem is equivalent to:

> Is there an infinite convex set (S\subset\mathbb Z) such that ((1_S\circ 1_S)(n)\to\infty) as (n\to\infty)?
> Or even ((1_S\circ 1_S)(n)\ge2) for all large (n)?

---
